\def\version{1}
\problemname{Home Delivery}
Carl has just moved away from home, and has realized that he now needs to buy food himself.
Since it is tiresome to go to the store, he instead orders the food online from the website \emph{Hemkör}, which delivers groceries directly to his door.

Carl is now buying food for the coming centuries.
In total, Carl has planned to eat $1 \le N \le 5\,000$ meals (numbered from $1$ to $N$) during the next $100\,000$ days.
Carl plans to eat the $i$th meal $1 \le P[i] \le 100\,000$ days from today, and it requires a total of $1 \le Q[i] \le 100$ kilograms of food.
Carl is not picky --- as long as he has $Q[i]$ kilograms of ingredients, it does not matter which ingredients he uses.

On Hemkör, there are $1 \le M \le 100\,000$ different grocery items for sale.
The $i$th item has a weight of $1 \le V[i] \le 100$ kilograms, a cost of $1 \le K[i] \le 2\,000$ kronor, and an expiration date that is $1 \le D[i] \le 100\,000$ days from today (an item can be used up \emph{to and including} the day it expires).
Carl can buy as many copies of each grocery item as he wants.

Carl has determined that it is possible to buy groceries so that he can prepare all meals.
Can you help him order groceries in a way that is also as cheap as possible?

\section*{Example}
Suppose Carl plans to prepare $2$ meals on days $1$ and $10$, with weights $4$ kilograms and $1$ kilogram.
There are three grocery items, with weights $4$, $2$, and $4$ kilograms, costs $10$, $2$, and $3$, and expiration dates $10$, $5$, $5$.
Carl can buy one copy of the first item and one copy of the second item.
For the first meal, he can use $2$ kilograms of each item.
For the second meal, he uses $1$ kilogram of the first item, which expires on the same day the meal is prepared.
One kilogram of the first item is left over and not used.

In total, this costs $10 + 2 = 12$ kronor.

\section*{Input}
The first line contains the positive integers $N$ and $M$.
Then follow $N$ lines, one for each meal.
The $i$th line contains the integers $P[i]$ and $Q[i]$.

Then follow $M$ lines, one for each item.
The $i$th line contains the integers $V[i]$, $K[i]$, and $D[i]$.

\section*{Output}
Output one number --- the minimum cost to buy food for all meals.

\section*{Scoring}
Your solution will be tested on a set of test groups. To earn points for a group, you must pass all test cases in that group.

\noindent
\begin{tabular}{| l | l | l |}
\hline
  Group & Points & Constraints \\ \hline
  1     & 31         & $N \le 100, M \le 100$ \\ \hline
  2     & 19         & $P[i] = 1$ for all meals \\ \hline
  3     & 15         & $V[i] = 1$ for all items \\ \hline
  4     & 35         & No additional constraints. \\ \hline
\end{tabular}
